\documentclass[10pt]{article}
\usepackage[T1]{fontenc}
\usepackage[utf8]{inputenc}
\usepackage[french]{babel}
\usepackage{amsmath,amssymb,mathrsfs,mathpazo}
\usepackage{geometry}
\geometry{a4paper,margin=2cm}
\usepackage{graphicx,array,tabularx,tabularray,colortbl,multirow,multicol,booktabs,enumitem,pifont,ulem,fancybox,diagbox,cancel}
\usepackage{pstricks,pstricks-add,pst-plot,pst-tree,pst-node,pst-text,pst-func,pst-math,pst-bspline,pst-3dplot,pst-eucl,pst-solides3d}
\let\pscalculate\undefined
\usepackage{pst-bezier}
\usepackage{tikz,pgfplots,tkz-tab}
\usetikzlibrary{arrows,arrows.meta,patterns,calc,positioning,trees,decorations.markings,shapes.geometric,shapes.misc,shapes.arrows,intersections,decorations.pathmorphing,decorations.pathreplacing,backgrounds,shadows,babel}
\pgfplotsset{compat=1.18}
\usepgfplotslibrary{fillbetween}
\usepackage[np]{numprint}
\usepackage{listings,makecell,pas-tableur,twoopt}
\let\barre\undefined
\usepackage[pstricks]{tablvar}
\definecolor{darkgreen}{RGB}{0,100,0}
\providecommand{\red}{\color{red}}
\providecommand{\blue}{\color{blue}}
\providecommand{\green}{\color{green}}
\providecommand{\black}{\color{black}}
\providecommand{\white}{\color{white}}
\providecommand{\gray}{\color{gray}}
\providecommand{\cyan}{\color{cyan}}
\providecommand{\magenta}{\color{magenta}}
\providecommand{\yellow}{\color{yellow}}
\providecommand{\violet}{\color{violet}}
\providecommand{\orange}{\color{orange}}
\providecommand{\euro}{\texteuro}
\providecommand{\EUR}{\texteuro}
\providecommand{\ieme}{\textsuperscript{e}}
\providecommand{\degres}{\ensuremath{^\circ}}
\providecommand{\degree}{\ensuremath{^\circ}}
\providecommand{\vv}[1]{\overrightarrow{#1}}
\providecommand{\Coord}[1]{#1}
\providecommand{\decosix}{}
\providecommand{\starredbullet}{\ensuremath{\bullet}}
\pagestyle{empty}
\setlength{\parindent}{0pt}
\usepackage{adjustbox}\DeclareUnicodeCharacter{2236}{:}\DeclareUnicodeCharacter{2190}{\ensuremath{\leftarrow}}
\providecommand{\R}{\mathbb{R}}
\providecommand{\N}{\mathbb{N}}
\providecommand{\Z}{\mathbb{Z}}
\providecommand{\Q}{\mathbb{Q}}
\providecommand{\D}{\mathbb{D}}
\providecommand{\C}{\mathbb{C}}
\providecommand{\e}{{\mathrm{e}}}
\providecommand{\vect}[1]{\overrightarrow{#1}}
\providecommand{\vectt}[1]{\overrightarrow{\mathrm{#1}}}
\providecommand{\vvt}[1]{\overrightarrow{\mathrm{#1}}}
\providecommand{\np}[2][]{{#2~#1}}
\providecommand{\vv}[1]{{\overrightarrow{#1}}}
\renewcommand{\d}{{\mathrm{d}}}
\providecommand{\degres}{{^\circ}}
\providecommand{\degre}{{^\circ}}
\providecommand{\blue}{\color{blue}}
\providecommand{\red}{\color{red}}
\providecommand{\enskip}{\quad}
\providecommand{\llbracket}{\lbrack\!\lbrack}
\providecommand{\rrbracket}{\rbrack\!\rbrack}
\providecommand{\euro}{{\text{€}}}
\providecommand{\no}{{\text{n°}}}
\providecommand{\interval}[2]{{[#1;#2]}}
\providecommand{\intervFO}[2]{{[#1;#2[}}
\providecommand{\dotfill}{\ldots}

\providecommand{\ds}{\displaystyle}
\providecommand{\textding}[1]{#1}
\providecommand{\Oij}{\(\mathrm{O};\vec{\imath},\vec{\jmath}\)}
\providecommand{\Oijk}{\(\mathrm{O};\vec{\imath},\vec{\jmath},\vec{k}\)}
\providecommand{\Ouv}{\(\mathrm{O};\vec{u},\vec{v}\)}
\providecommand{\intervFF}[2]{\ensuremath{[#1;#2]}}
\renewcommand{\euro}{\ifmmode\text{\texteuro}\else\texteuro\fi}
\DeclareUnicodeCharacter{00B0}{\ifmmode^{\circ}\else\textdegree\fi}
\DeclareUnicodeCharacter{20AC}{\ifmmode\text{\texteuro}\else\texteuro\fi}
% Énoncés : Éducation nationale. Transcription et corrigés : APMEP, Denis Vergès et contributeurs. Découpage et classement : CoopMaths, Sébastien Lozano, Éric Elter et contributeurs. Adaptation de la présentation : Maths Explorer.
% CC BY-SA : https://coopmaths.fr/www/about/
% https://coopmaths.fr/alea/static/e3c/2020/tex/e3c_2020_00_sujet32_3.tex
\begin{document}
\noindent\textit{APMEP / CoopMaths. CC BY-SA. Présentation adaptée par Maths Explorer.}\par\medskip
\medskip

Dans le repère ci-dessous, on note $\mathcal{C}_f$ la courbe représentative d'une fonction $f$ définie sur l'intervalle $[-10~;~2]$.

On a placé dans ce repère les points $A~(0~;~2$), $B~(2~;~0)$ et $C~(-2~;~0)$.

On dispose des renseignements suivants :
\begin{itemize}
\item Le point B appartient à la courbe $\mathcal{C}_f$.
\item La droite (AC) est tangente en A à la courbe $\mathcal{C}_f$.
\item La tangente à la courbe $\mathcal{C}_f$ au point d'abscisse 1 est une droite parallèle à l'axe des abscisses.
\end{itemize}


\psset{unit=1.05cm,labelFontSize=\scriptstyle, showorigin=false}
\begin{pspicture}(-10.35,-1.3)(2.7,3)
\multido{\n=-10.33333333+0.33333333}{40}{\psline[linewidth=0.3pt,linecolor=lightgray](\n,-1.33)(\n,3)}
 \multido{\n=-1.3333333+0.3333333}{14}{\psline[linewidth=0.3pt,linecolor=lightgray](-10.4,\n)(2.7,\n)}

 
 \psaxes[linewidth=1.25pt]{->}(0,0)(-10.33,-1.33)(2.7,3)
\def\Func{2.71828 x exp 2 x sub mul}
 \psplot[plotpoints=3000,linewidth=1.25pt,linecolor=green]{-10}{2}{\Func}
 \psplot[plotpoints=3000,linewidth=0.85pt,linecolor=red]{-10.333}{2.3}{2.71828}
 \psplot[plotpoints=3000,linewidth=0.85pt,linecolor=blue]{-3.4}{1.02}{x 2 add}
 \psdots[dotstyle=Mul,dotscale=2.6](0,2)(2,0)(-2,0)
 \uput[ur](2,0){B} \uput[r](0,2){A} \uput[u](-2,0){C} \uput[dl](0,0){O} 
 \uput[dr](1.8,1.32){\blue $\displaystyle \mathcal{C}_f$}
\end{pspicture}

\begin{enumerate}
\item Déterminer la valeur de $f'(1)$.
\item Donner une équation de la tangente à la courbe $\mathcal{C}_f$ au point A.
\end{enumerate}
On admet que cette fonction $f$ est définie sur $[-10~;~2]$ par \[f(x)=(2-x)\e^x.\]
\begin{enumerate}[resume]
\item Montrer que pour tout réel $x$ appartenant à l'intervalle $[-10~;~2]$,

\[f'(x)=(-x+1)e^x.\]

\item En déduire le tableau de variations de la fonction $f$ sur l'intervalle $[-10 ;2]$.
\item Déterminer une équation de la tangente à la courbe $\mathcal{C}_f$ au point B.
\end{enumerate}

\vspace{0,5cm}

\end{document}
